% LaTeX transcription of the original AI-written research note, 18 November 2026.
% The original PDF is preserved unchanged alongside this file.
% The original .tex was not available in the linked chat or Library; this
% compilable transcription was prepared on 19 November 2026 from that PDF.
% Build: pdflatex polylog_sorting_adversary.tex (run twice for links).
\documentclass[11pt,a4paper]{article}
\usepackage[margin=26mm,top=22mm,bottom=22mm]{geometry}
\usepackage{amsmath,amssymb,microtype}
\usepackage[hidelinks]{hyperref}
\hypersetup{
  pdftitle={An optimal sorting adversary in polylogarithmic time per comparison},
  pdfsubject={A word-RAM refinement of the order-restriction DAG adversary}
}
\setlength{\parindent}{0pt}
\setlength{\parskip}{5pt}
\begin{document}
\begin{center}
{\large\bfseries An optimal sorting adversary in polylogarithmic\\
time per comparison\par}
\vspace{6pt}
Research note\quad 18 November 2026
\end{center}
\vspace{6pt}

\textbf{Theorem.} On a deterministic word RAM with $\Theta(\log n)$-bit words,
one uniform online adversary for $n$ distinct labels answers every comparison in
$O((\log n)^C)$ worst-case time for an absolute constant $C$. Every correct
comparison sorter uses at least $(1-o(1))n\log_2 n$ comparisons against it.
Initialization uses $O(n)$ time, and the space bound is $n\,\mathrm{polylog}\,n$ words.

The starting point is Nash's order-restriction DAG adversary [1], version 2.
Its local calculations for fixed accuracy inspect a bounded number of vertices,
but their inputs can have many fractional bits. Two observations yield the
stronger time bound. An absolute approximation suffices for the probability
certificates. When a certificate fails, the only new number required is a point
well inside an interval with exact endpoints. That point can be constructed by
copying a binary prefix and adding a small multiple of a power of two.

\section{What the comparison proof requires}
All logarithms are binary. Fix $k\ge64$ and put $H=2^k$ and
$\eta=1/(100k)$. Each node of the adversary's tree stores a DAG of free labels
and an interval $(\ell_x,r_x)$ for every free label $x$. A feasible assignment
chooses a point in each interval, increasing along each arc. The intervals are
canonical, meaning that their endpoints are the infimum and supremum attainable
by the coordinate. A child at position $c$ represents a block of labels. A
refinement fixes $x$ at $c$ in the parent, deletes its coordinate there, and gives
it a fresh unconstrained coordinate in that child. Let $f_x$ denote the marginal
density of $x$ for a uniform feasible assignment in the parent.

The proof of Lemma 3.6 and Theorem 3.7 of [1] works for any deterministic choices
satisfying the following three estimates.
\begin{itemize}
\item A certified answer retains at least $1/2-\eta$ of the current node volume.
\item A refinement requested by \textsc{Select} has $\log f_x(c)>k$.
\item Every refinement, including a repair, selects an interior dyadic $c$.
In reduced form $c=(2i+1)/2^{m+1}$, it satisfies
$\log f_x(c)\ge(1-1/k)m-4$.
\end{itemize}
We keep the path and degree repair rules of [1]. Under these hypotheses, if $T_q$
is the number of total orders consistent with the first $q$ answers, that proof gives
\begin{equation}
\log T_q\ge\log(n!)-(1+\delta_k)q,\qquad \delta_k=O(\log k/k).
\end{equation}
In particular, allowed choices can change which answer the adversary returns.
The proof uses the three inequalities and the legality of interior refinements.

\newpage
\section{Short inputs suffice for the local certificates}
Theorem B.1 of [1] supplies a finite list of intervals enclosing each queried
coordinate. A uniform entry can be coupled with a uniform feasible point so that
the point lies inside its intervals. Except for an $\alpha$ fraction of entries,
the intervals have width at most $\delta$. The tolerances $\alpha,\delta>0$ used
in Appendix D of [1] depend only on $k$. The number of local vertices, digit
assignments and arithmetic steps also depends only on $k$.

Run the enclosure with width tolerance $\delta/2$ and the same $\alpha$. Fix a
dyadic grid of mesh $u\le\delta/4$. In the lower computation round every input
down to this grid, including both interval walls at every vertex and the initial
positions. In the upper computation round every input up. Carry out all
subsequent updates exactly, with the original digit sequence for each list entry.

To check the effect on one trajectory, each update takes maxima and minima and
then a convex combination. These operations are monotone and nonexpansive for
the maximum norm, jointly in the state and the stored walls. By induction, each
computed lower endpoint lies between the original lower endpoint minus $u$ and
the original endpoint. The upper endpoint increases by at most $u$. Thus the
wider intervals still enclose the same feasible point under the original
coupling, and every formerly good entry has width at most $\delta/2+2u\le\delta$.
Both conclusions of Theorem B.1 hold at the tolerances required in Appendix D.
For a comparison against an existing anchor, the resulting short endpoints can
be compared with that exact anchor.

The resulting list can therefore be used in the original mass tests. A
successful comparison certifies its retained volume. A successful central-interval
test supplies a short-dyadic interval $J$ with mass at least $1/4$ in $J$ and on
each side of it. Applying the original rounding rule in $J$ gives
$\log f_x(c)\ge m-4$ by Lemmas A.6 and E.1 of [1]. All temporary arithmetic in
the sweeps and in this successful anchor calculation has bit length bounded by a
computable function of $k$.

\section{An exact anchor when certification fails}
Put $B=(k+1)^{2H+1}$, the bound on the number of coordinates comparable with one
free coordinate. A failed test implies that the chosen coordinate's canonical
width $W=r-\ell$ satisfies $W<2^{-2kB}$, by Lemma D.1 of [1]. Write
$L=\log(1/W)$. Throughout the middle half of $(\ell,r)$ the density bound (18)
of [1] gives
\[
\log f_x(c)\ge L-2B>(1-1/k)L,\qquad L-2B>k.
\]
Choose the largest power of two $h\le W/4$, and define
\begin{equation}
c=h\lfloor\ell/h\rfloor+3h.
\end{equation}
Since $4h\le W<8h$, we have $2h<c-\ell\le3h$, which puts $c$ in the middle
half. Since $c$ is a multiple of $h$ and $W<8h$, its reduced denominator
exponent satisfies $m<L+2$. This yields $\log f_x(c)\ge(1-1/k)m-4$, as required.
All three estimates on page 1 follow with the original constants. The formula
requires only the scale of the width, a prefix of the lower endpoint and two
power-of-two additions.

\newpage
\section{Working with long dyadics}
Before the horizon $Q=n\lceil\log n\rceil$, every query moves at most two queried
labels before the answer. A moved label has density one in its fresh coordinate
and cannot be selected again during that loop. After the answer, the single
repair moves at most $H$ labels (Algorithm 7 of [1]). Let $A(k)$ bound the
precision of anchors from successful central tests. If the stored endpoints
have at most $M$ fractional bits, then $W\ge2^{-M}$, so (2) uses at most $M+2$
bits. A sufficient common capacity for the number of fractional bits is therefore
\[
D=A(k)+2(Q+1)(H+2)+8.
\]
We shall choose $A(k),H\le\log n$ for all sufficiently large $n$. In particular,
$D$ is polynomial in $n$, though a single endpoint can be much longer than one word.

Represent a binary number, including a constant number of integer bits, by a
persistent complete binary tree on its bit positions. Pad its capacity to a
power of two. Equal blocks at a given height share one identifier. A
deterministic balanced search tree maps each triple (height, left identifier,
right identifier) to its unique identifier. Store one all-zero and one all-one
block of each height. Thus the initial tree uses $O(\log D)$ records. New
numbers reuse every unaltered subtree. The following operations touch
$O(\log D)$ blocks each, with a dictionary lookup for every new block.
\begin{itemize}
\item Compare two numbers by descending to the first unequal bit, skipping
equal child identifiers.
\item Truncate after a specified bit by clearing a suffix. To add one unit at
a specified bit, locate the rightmost preceding zero at or before that bit,
set it, and clear the intervening run of ones. Uniform blocks allow the search
and both range updates in logarithmic time.
\end{itemize}
Find the least $s\ge0$ for which $\ell+2^{-s}\le r$ by binary search on bit
positions. Then $h=2^{-(s+2)}$ has the defining property in (2). Truncate $\ell$
to this grid and add $3h$ using two power-of-two additions. This costs
$O(\log D)$ tree operations. For rounding up, a comparison with the truncated
value detects a nonzero tail. Reading the bounded prefix for a local sweep,
or embedding a short computed endpoint for exact comparison, costs a computable
function of $k$ times a power of $\log D$. Throughout, equality of identifiers
is exact.

\section{Graph maintenance in worst-case time}
Between calls the node DAGs have degree below $k$ and directed paths of fewer
than $H$ vertices. Immediately after an arc insertion the bounds are $k$ and
fewer than $2H$. Thus all the ancestors and descendants of a changed vertex
occupy a set of size bounded in terms of $k$. Inserting $a\to b$ propagates
$\ell_a$ only to $b$ and its descendants and $r_b$ only to $a$ and its ancestors.
A one-sided interval bound propagates in one direction. Fixing $x$ at $c$
changes only the upper bounds of its predecessors and the lower bounds of its
successors, by formula (14) of [1]. All these changes use exact comparisons
between endpoints.

\newpage
Every newly forbidden path must contain the inserted arc, and only its
endpoints can newly attain degree $k$. Searching the two cones of that arc
finds a path for repair, including the lexicographically first one if required.
The at most $H$ moves in such a repair take bounded local graph work per move
for fixed $k$. The graph visited by the enclosure has bounded undirected
radius and bounded degree as functions of $k$, and can be collected from
adjacency lists.

Store each label's current tree node. The hierarchy grows by leaf additions.
Pointers to ancestors at power-of-two distances support insertion of a node,
lowest-common-ancestor queries, and finding the first child below that ancestor
in $O(\log N)$ operations for $N$ tree nodes. This locates the node at which a
comparison is resolved. Each parent's children can be stored in a balanced
dictionary of exact dyadic anchors. There are at most $1+(H+2)Q$ tree nodes
before the horizon.

\paragraph{After the horizon.}
Freeze the state at $Q$ and use the following fixed order. At each node place
an anchor at its own value and each free item $x$ immediately to the right of
$\ell_x$. Break ties among items with equal lower bounds by the length of the
longest directed path ending at $x$, then by label. A finite node admits a
sufficiently small positive perturbation realizing these choices. Use ordered
offsets smaller than any positive gap between distinct endpoints or anchors.
Every item is below its upper bound; on an arc, the lower bounds are
nondecreasing and the path length increases. Thus these local orders give a
represented global order, which respects every previous comparison. A query in
the frozen state needs only the lowest common ancestor and the predecessor
cones of at most two items. Switching modes requires one flag and works even
if comparisons continue indefinitely.

\paragraph{One uniform algorithm.}
For fixed $k$, the foregoing implementation has worst-case query time
$F(k)(\log n)^c$ for a fixed $c$ and a computable $F$. Here $F$ includes
arithmetic on short endpoints, all finite enumerations and graph searches,
the repair batch, and bounds for $A(k)$ and $H$. All loops involving the local
parameters have computable bounds from the construction. The variable-length
operations use a fixed number of nested logarithmic searches. The interning
dictionary gets at most a polynomial in $n$ many entries when $F(k)\le\log n$,
since each of $Q$ calls creates at most $F(k)(\log D)^{O(1)}$ blocks,
independently of dictionary lookup time. Its identifiers and addresses therefore
fit in $O(\log n)$ bits.

To choose $k$, simulate for $\lfloor\log n\rfloor$ steps a fixed program emitting
successive complete parameter tuples and upper bounds $F(64),F(65),\ldots$.
Use the largest completed stage with $F(k)\le\log n$. Every fixed stage
eventually qualifies, so $k(n)\to\infty$. Use the label order for the finitely
many sizes with no qualifying stage. Initialization can create the $n$ isolated
label records in linear time. The per-query time is now polylogarithmic, and
space through the horizon is $n\,\mathrm{polylog}\,n$. A sorter stopping before
$Q$ has $T_q=1$, so (1) forces $q\ge\log(n!)/(1+\delta_{k(n)})$. A sorter
reaching $Q$ already meets the claimed bound. This proves the theorem.
\hfill$\square$

The convergence of the coefficient is governed by the cost of the finite
enumeration as a function of $k$. The implementation above gives an asymptotic
existence result.

\vfill
{\small\textbf{Reference.} [1] Nicholas Nash, \emph{An Optimal Deterministic
Adversary for Sorting}, version 2, November 2026.
\href{https://hexagonmath.org/2611.00196v2}{hexagon:2611.00196v2}.}
\end{document}
